\documentclass[11pt,a4paper]{article}

\usepackage[spanish]{babel}
\usepackage[utf8]{inputenc}
\usepackage[T1]{fontenc}
\usepackage{lmodern}
\usepackage{microtype}
\usepackage{geometry}

\geometry{margin=2.8cm}

\title{\textbf{De lo propio, de lo común y de su igualdad}}
\author{Vicente Domínguez Arca}
\date{}

\begin{document}

\maketitle

En física hay un recurso metafórico utilizado con enorme frecuencia: la escala. Etimológicamente, la palabra hace referencia a un aumento progresivo y ordenado, como una sucesión de escalones. La metáfora resulta clara: la escala es la regla con la que decidimos medir un fenómeno.

Toda regla puede estar graduada de distintas maneras. Puede medir milímetros, centímetros o metros, y no por ello deja de ser la misma regla. En física ocurre algo semejante. Una misma realidad puede observarse desde distintas escalas, y con cada cambio algunas leyes se vuelven dominantes mientras otras pasan a ser prácticamente despreciables.

En nuestra escala humana, por ejemplo, la gravedad gobierna buena parte de cuanto ocurre. Sin embargo, cuando descendemos a escalas micrométricas o nanométricas, la gravedad pierde casi toda su importancia frente a fenómenos como la electrostática o la electrodinámica. La realidad no cambia; cambia aquello que resulta relevante para describirla.

No ha sido un trabajo sencillo determinar dónde se producen esas transiciones. Buena parte de la física moderna ha dedicado esfuerzos monumentales a comprender en qué punto una determinada descripción deja de ser suficiente y otra comienza a imponerse.

Sin embargo, existen fenómenos que parecen escapar a ese cambio. Son aquellos que presentan lo que se conoce como \textbf{invarianza de escala}. Es decir, independientemente del nivel desde el que los observemos, conservan esencialmente la misma estructura.

Un ejemplo sencillo puede encontrarse en la geometría. Existen figuras cuyo borde mantiene siempre el mismo perfil por mucho que ampliemos una pequeña porción de él. Son los conocidos fractales. Hay toda una familia de teorías alrededor de estos objetos, pero aquí únicamente nos interesa una idea: algunas formas permanecen siendo ellas mismas aunque cambiemos la escala desde la que las observamos.

Hasta aquí podría terminar este texto. Sospecho que a algún lector atento le habrá bastado el título y esta breve introducción para anticipar hacia dónde quiero dirigirme. Sin embargo, las intuiciones también necesitan tiempo para desplegarse. Siempre quedan embrollos que aclarar, muebles que limpiar y suelos que fregar.

Definir etimológicamente lo propio ha ocupado a filólogos e historiadores durante siglos. Aún hoy existe cierta controversia sobre su origen último, aunque puede alcanzarse un consenso razonable: lo propio es aquello que pertenece a un sistema individual.

Conviene detenerse un instante en esa última expresión. Hablo de sistema individual, no necesariamente de individuo. Un sistema es cualquier realidad que podemos aislar de un conjunto mayor, ya sea mediante fronteras materiales o mediante límites conceptuales. No hace falta resolver aquí una cuestión ontológica; basta con aceptar que somos capaces de distinguir unas cosas de otras para poder hablar de ellas.

Desde esa perspectiva, parece sencillo reconocer qué es lo propio.

Es propio del árbol el lugar que ocupa.

Es propio del automóvil la plaza donde permanece estacionado.

Son propios de una persona los pantalones que viste.

Son propias sus gafas.

Hasta aquí todo parece evidente.

Pero la evidencia comienza a resquebrajarse en cuanto cambiamos la escala.

Es propio de la raíz el recorrido que describe bajo la tierra.

Es propio del neumático el fragmento de suelo sobre el que descansa.

Son propios de los pantalones los tobillos que alcanzan.

Son propias de las gafas las orejas sobre las que reposan sus patillas.

Lo propio, entonces, parece sufrir también una forma de invarianza de escala. No depende tanto del objeto como del sistema que decidimos considerar. Allí donde definimos un sistema aparece aquello que le es propio. Basta modificar la escala para que cambie el conjunto de relaciones que observamos, sin que desaparezca la noción misma de pertenencia.

Demos ahora un paso más. Y digo ``demos'' porque confío en haber conseguido algo difícil: que lector y escritor formemos durante unos instantes un mismo sistema.

Yo escribo un texto que he pensado.

El lector lee ese texto, lo piensa también, quizá lo critica, quizá lo continúa.

Hay algo que nos une.

Hay algo común.

Pero ¿qué es, exactamente, lo común?

La etimología suele asociarlo con aquello que reúne deberes, beneficios, garantías, servicios, lugares, dones o materias entre varios. Lo común es aquello que sostiene simultáneamente a más de un sistema.

Es común la tierra de un monte para los árboles que crecen sobre ella.

Es común una zona de aparcamiento para los automóviles que la ocupan.

Son comunes unas piernas para todos los estilos de pantalón que puedan vestirlas.

Son comunes unas orejas para todas las gafas que puedan descansar sobre ellas.

Y entonces aparece una sospecha.

Quizá la diferencia entre lo propio y lo común sea menos profunda de lo que hemos creído. Quizá sea únicamente una diferencia de escala.

Porque, si nos atenemos a la etimología, aquello que llamamos común no deja de comportarse como lo propio cuando ampliamos el sistema de referencia. Lo que era propio de una relación pasa a ser común respecto de un conjunto mayor de relaciones. Nada parece haber cambiado salvo el tamaño del sistema que hemos decidido observar.

Esta posibilidad me resulta especialmente sugerente porque vivimos en una época que invierte enormes esfuerzos ---desde la política, la economía, la administración o la sociología--- en separar cuidadosamente lo propio de lo común. Gran parte del pensamiento contemporáneo, especialmente bajo las lógicas neoliberales y capitalistas, necesita que esa frontera permanezca perfectamente delimitada. Allí donde existen límites aparecen gradientes; donde hay gradientes aparecen flujos; y donde existen flujos puede circular el capital.

No pretendo simplificar problemas cuya complejidad excede con mucho estas páginas. Tan solo me pregunto si parte del llamado problema de lo común nace de haber aceptado como real una distinción que quizá solo pertenezca a nuestra descripción.

Quizá lo común sea, sencillamente, lo propio observado desde otra escala.

La cuestión que realmente quisiera dejar abierta no es si esta idea resulta convincente para quien la lea.

La cuestión es otra.

¿Y si fuera verosímil?

¿Y si hubiéramos vivido demasiado tiempo sometidos a una imagen que ha organizado nuestras instituciones, nuestros conflictos y nuestras formas de entender la propiedad sin corresponderse realmente con la naturaleza de las relaciones?

¿Qué ocurriría si aceptáramos que lo común no es lo contrario de lo propio, sino únicamente lo propio contemplado desde otra escala?

\end{document}